\chapter{Termo de Abertura do Projeto}

\section{Dados do projeto}
\begin{description}
    \item [Nome do Projeto:] Micromouse - PI1
    \item [Data de abertura:] 16/03/2026
    \item [Código:] FGA0303 (Grupo 1 - Juliana)
    \item [Patrocinador:] Universidade de Brasília
    \item [Gerente:] André Cláudio Maia da Cunha
        \begin{description}
            \item [Matrícula:] 222037648
            \item [E-mail:] 222037648@aluno.unb.br
            \item [Telefone:] (61) 99663-7828
        \end{description}
\end{description}

\section{Objetivos}

O projeto tem como principal objetivo desenvolver e construir um robô autônomo com chassi em impressão 3D, capaz de resolver labirintos desconhecidos de dimensões 4x4, 8x8 e 16x16 células. O sistema integrará componentes mecânicos, de hardware embarcado
e algoritmos de exploração, como o Flood Fill, para detectar paredes, realizar mapeamento incremental e alcançar o centro do trajeto no menor tempo possível sem qualquer tipo de intervenção humana. Além do robô, o projeto engloba a construção de um labirinto
modular de MDF padronizado para testes e um sistema web de telemetria para visualização de dados de navegação e energia em tempo real.

Para avaliar a eficácia do projeto e atestar a integração multidisciplinar técnica do produto desenvolvido, foram estabelecidos os seguintes indicadores-chave de desempenho (KPIs), os quais irão nortear as entregas acadêmicas do projeto:

\begin{enumerate}
    \item \textbf{Autonomia e Resolução:} Capacidade do robô de alcançar o objetivo central dos labirintos sem intervenções humanas, mapeando incrementalmente as células usando seus módulos sensoriais.
    \item \textbf{Eficiência do Algoritmo de Navegação:} Tempo e quantidade de tentativas necessárias para a resolução da rota, atestando o funcionamento de lógica embarcada.
    \item \textbf{Precisão da Telemetria:} Confiabilidade da comunicação sem fio do ESP32, garantindo o envio, processamento e visualização em tempo real de informações como posição, tempo, velocidade média e bateria com baixa latência. 
    \item \textbf{Resistência, Escalabilidade e Peso:} Integridade do chassi e obediência à restrição de tamanho máximo de 25 cm x 25 cm para largura e comprimento, de modo a assegurar que não haja danos no contato com as paredes. 
    \item \textbf{Autonomia Energética:} Duração operacional contínua mínima de 30 minutos em pista, validando o circuito \textit{step-down} construído e o correto dimensionamento das baterias de NiMH mediante margem de segurança contra picos de corrente.
\end{enumerate}

\section{Mercado-alvo}

O mercado-alvo deste projeto abrange a comunidade acadêmica e entusiastas das engenharias e robótica competitiva. Trata-se de um sistema direcionado inicialmente a estudantes, professores, pesquisadores das áreas de eletrônica, software, energia, estruturas e participantes de feiras tecnológicas e eventos como competições micromouse.
Além disso, o projeto possui potencial expressivo de uso como plataforma didática para o ensino em disciplinas de engenharia integrando hardware, software, sensoriamento e controle dinâmico de estruturas físicas.

\section{Requisitos}

Os requisitos foram analisados e revisados a partir da priorização de itens indispensáveis ao sucesso do Micromouse, de acordo com o escopo e as regras da competição simulada.

\begin{landscape}
    % Força a tabela a ignorar recuos laterais e ocupar o centro/largura total
    \setlength{\LTleft}{0pt}
    \setlength{\LTright}{0pt}

    % O @{\extracolsep{\fill}} faz com que a tabela se estique até a margem
    \begin{longtable}{@{\extracolsep{\fill}}|l|p{14.5cm}|l|l|}
        
        \caption{Requisitos Funcionais e Não Funcionais} \label{tab:requisitos} \\
        
        \hline
        \textbf{ID} & \textbf{Requisito} & \textbf{Categoria} & \textbf{Prioridade} \\ \hline
        \endfirsthead

        % Cabeçalho repetido (sem texto "continuação")
        \hline
        \textbf{ID} & \textbf{Requisito} & \textbf{Categoria} & \textbf{Prioridade} \\ \hline
        \endhead

        % Rodapé simples (apenas a linha, sem texto "continua")
        \hline
        \endfoot

        \hline
        \endlastfoot

        % --- REQUISITOS FUNCIONAIS ---
        \multicolumn{4}{|c|}{\textbf{REQUISITOS FUNCIONAIS (RF)}} \\ \hline
        RF01 & O sistema deve navegar de forma autônoma em labirintos 4x4, 8x8 e 16x16 após a ativação. & Navegação & Must have \\ \hline
        RF02 & O sistema deve detectar paredes e mapear o ambiente progressivamente em tempo real. & Navegação & Must have \\ \hline
        RF03 & O sistema deve decidir a rota utilizando algoritmos de resolução de labirintos (ex: Flood Fill). & Navegação & Must have \\ \hline
        RF04 & O sistema deve identificar o momento em que o robô atinge o centro do labirinto. & Navegação & Must have \\ \hline
        RF05 & O sistema deve calcular e percorrer a rota mais rápida de retorno à origem após o mapeamento. & Navegação & Could have \\ \hline
        RF06 & O sistema deve processar dados dos sensores de distância continuamente para detectar obstáculos e paredes. & Sensoriamento & Must have \\ \hline
        RF07 & O sistema deve monitorar sua posição continuamente via odometria e controle de orientação. & Sensoriamento & Must have \\ \hline
        RF08 & O sistema deve aferir o nível de bateria via leitura analógica com divisor de tensão. & Sensoriamento & Must have \\ \hline
        RF09 & O sistema deve transmitir dados de telemetria sem fio para o servidor. & Transmissão de Dados & Must have \\ \hline
        RF10 & O sistema web deve exibir telemetria em tempo real (tipo de labirinto, velocidade, trajeto, tempo e bateria). & Visualização & Must have \\ \hline
        RF11 & O sistema web deve exibir um indicador de status atestando o cumprimento da missão. & Visualização & Must have \\ \hline
        RF12 & O sistema deve armazenar dados de corridas concluídas de forma persistente em banco de dados. & Armazenamento & Must have \\ \hline
        RF13 & O sistema web deve permitir consulta ao histórico e comparação de desempenho entre execuções. & Visualização & Should have \\ \hline
        RF14 & O sistema deve identificar as dimensões do labirinto automaticamente ou via seletor físico antes da ativação. & Navegação & Must have \\ \hline
        RF15 & O sistema web deve permitir filtrar o histórico de dados por dimensões do labirinto ou exibir todos. & Visualização & Should have \\ \hline

        % --- REQUISITOS NÃO FUNCIONAIS ---
        \multicolumn{4}{|c|}{\textbf{REQUISITOS NÃO FUNCIONAIS (RNF)}} \\ \hline
        RNF01 & O robô deve possuir no máximo 25 cm de comprimento e de largura. & Restrição Física & Must have \\ \hline
        RNF02 & O robô não deve usar motores de combustão, voar, pular ou danificar o labirinto. & Restrição Física & Must have \\ \hline
        RNF03 & O banco de baterias deve representar no máximo 25\% do peso total do robô. & Restrição Física & Should have \\ \hline
        RNF04 & O processamento do sistema embarcado deve utilizar exclusivamente o microcontrolador ESP32. & Tecnologia & Must have \\ \hline
        RNF05 & O sistema de alimentação deve prover autonomia mínima de 30 minutos contínuos. & Energia & Must have \\ \hline
        RNF06 & O sistema deve computar e decidir a rota na célula atual em até 5 segundos. & Desempenho & Must have \\ \hline
        RNF07 & A interface web deve exibir telemetria com latência máxima de 500 ms. & Desempenho & Must have \\ \hline
        RNF08 & Em parada de emergência, os atuadores devem cessar movimento em no máximo 2 segundos. & Confiabilidade e Segurança & Must have \\ \hline
        RNF09 & O código-fonte e o mapa não devem sofrer modificações remotas durante a execução na pista. & Segurança e Confiabilidade & Must have \\ \hline
        RNF10 & O banco de dados para armazenamento de corridas deve possuir arquitetura relacional. & Tecnologia & Should have \\ \hline

    \end{longtable}
\end{landscape}

\section{Justificativa}

O desenvolvimento do Micromouse configura-se como um campo de testes real e prático das disciplinas da Faculdade de Ciências e Tecnologias em Engenharia (FCTE), exigindo um corpo de estudantes não apenas o conhecimento teórico, mas habilidades de execução. Assim, a justificativa deste projeto
se fundamenta na necessidade inerente de preparar futuros profissionais em sistemas interligados de áreas multidisciplinares de hardware, software e estrutural aplicados ao mundo real.

\section{Indicadores}

A listagem a seguir explicita até dez indicadores numéricos que evidenciam o impacto mercadológico e aplicabilidade prática do sistema como objeto técnico de robótica educacional e competitiva:

\begin{enumerate}
    \item Percentual de precisão na detecção e mapeamento de paredes do labirinto.
    \item Número estimado de laboratórios no DF que poderiam usar o modelo para ensino de eletrônica e programação.
    \item Tempo médio (em milissegundos) de comunicação entre o sistema web de telemetria e o microcontrolador ESP32 na pista.
    \item Custo total do chassi 3D (PLA) e módulos embarcados comparado a kits proprietários tradicionais do mercado.
    \item Autonomia média suportada pelo banco de baterias dimensionado frente às atividades dos motores e placas.
    \item Quantidade de commits e issues de projetos gerados no Github atestando o engajamento dos alunos nas Sprints propostas.
    \item Índices de sucesso da navegação autônoma em pistas progressivamente mais complexas (4x4 vs 16x16).
\end{enumerate}